\begin{helper}
Let \(D\) be a digraph on a finite vertex set \(W\), let \(3\le\ell\), let
\(4\le s\le n\), and put
\[
  \Omega=\{\phi:\ [\ell-1]\to([n]\to W)\}.
\]
Assume
\[
  |W|>2^{2s-4},\qquad
  \overrightarrow{i_s}(D)\le 2^{3s^2/2},
\]
and
\[
  n\le 2^{(s/2-4)(\ell-1)}.
\]
If \(\phi\in\Omega\) is called bad when
\(\noderef{RandomHomomorphismColoring}(D,\ell,n,\phi)\) has a final-color
monochromatic \(s\)-clique, then the number of bad \(\phi\) is strictly smaller
than \(|\Omega|\).
\end{helper}
\begin{proof}
By \(\noderef{RandomHomomorphismFinalColorBadCountUpper}\), the number of bad
\(\phi\) is at most
\[
  \binom ns\,\overrightarrow{i_s}(D)^{\ell-1}\,
  |W|^{(\ell-1)(n-s)}.
\]
The total number of map families is
\[
  |\Omega|=|W|^{(\ell-1)n}.
\]
Since \(\binom ns\le n^s\), it is enough to compare
\[
  n^s\,\overrightarrow{i_s}(D)^{\ell-1}
\]
with \(|W|^{s(\ell-1)}\).

The hypotheses give
\[
  \overrightarrow{i_s}(D)^{\ell-1}
  \le 2^{(3s^2/2)(\ell-1)}.
\]
Also \(|W|>2^{2s-4}\), and since \(s(\ell-1)>0\), raising to the
\(s(\ell-1)\)-st power gives the strict inequality
\[
  2^{(2s-4)s(\ell-1)}<|W|^{s(\ell-1)}.
\]
Multiplying this strict comparison by the positive factor
\[
  2^{(-s^2/2+4s)(\ell-1)}
\]
and using the exponent identity
\[
  (2s-4)s(\ell-1)+(-s^2/2+4s)(\ell-1)
  =(3s^2/2)(\ell-1),
\]
we obtain
\[
  \overrightarrow{i_s}(D)^{\ell-1}
  < |W|^{s(\ell-1)}\,2^{(-s^2/2+4s)(\ell-1)}.
\]
Finally, \(\noderef{RandomHomomorphismFinalColorArithmetic}\) applied to the
assumed bound on \(n\) gives
\[
  n^s\,2^{(-s^2/2+4s)(\ell-1)}\le 1.
\]
Therefore
\[
  n^s\,\overrightarrow{i_s}(D)^{\ell-1}<|W|^{s(\ell-1)}.
\]
Multiplying by the positive factor \(|W|^{(\ell-1)(n-s)}\) and using
\[
  s(\ell-1)+(\ell-1)(n-s)=(\ell-1)n
\]
shows that the upper bound for bad maps is strictly smaller than
\(|W|^{(\ell-1)n}=|\Omega|\).  Hence the bad subtype has cardinality strictly
less than \(\Omega\), as claimed.
\end{proof}

